% =====================================================================
%  tikz-protypa.tex — τα ΠΡΟΤΥΠΑ ΣΧΗΜΑΤΑ της εργασίας
%
%  ΟΙΚΟΓΕΝΕΙΑ Α — ΠΟΣΟΤΙΚΑ (άξονες, αριθμοί).  Μόνο pgfplots.
%  ΟΙΚΟΓΕΝΕΙΑ Β — ΔΟΜΙΚΑ (ροή, καταστάσεις).  TikZ, με πηγή σε Mermaid.
%
%  Χρησιμοποιεί ΜΟΝΟ ό,τι φορτώνει ήδη το preamble.tex.
%  ΟΛΑ ΤΑ ΜΠΛΟΚ ΜΕΤΑΓΛΩΤΤΙΣΤΗΚΑΝ ΜΕ XeLaTeX ΠΡΙΝ ΜΠΟΥΝ ΕΔΩ.
% =====================================================================

% ---------------------------------------------------------------------
%  ΚΟΙΝΟ ΥΦΟΣ — αντέγραψέ το ΜΙΑ ΦΟΡΑ στο preamble.tex και μετά κάθε
%  σχήμα λέει μόνο τα δεδομένα του. Αυτό είναι το «τυποποιημένο».
% ---------------------------------------------------------------------
%
% \pgfplotsset{
%   thesisaxis/.style={
%     width=0.80\linewidth, height=5.6cm,
%     grid=major, grid style={cGrey!20, line width=0.3pt},
%     axis line style={cGrey!70, line width=0.5pt},
%     tick label style={font=\scriptsize, /pgf/number format/1000 sep={}},
%     label style={font=\small},
%     legend style={at={(0.5,-0.28)}, anchor=north, legend columns=2,
%                   font=\scriptsize, draw=cGrey!40, fill=white,
%                   /tikz/every even column/.append style={column sep=8pt}},
%     every axis plot/.append style={thick, mark size=2pt},
%   },
%   thesisbar/.style={thesisaxis, ybar, bar width=7pt, ymin=0,
%     nodes near coords, nodes near coords style={font=\tiny, text=cGrey},
%     enlarge x limits=0.12},
% }
%
% ΣΕΙΡΑ ΧΡΩΜΑΤΩΝ, ΣΤΑΘΕΡΗ, ΠΟΤΕ ΚΥΚΛΙΚΗ:
%   1. cCFT   #0072B2   2. cBFT  #D55E00   3. cAux #009E73
%   4. cSand  #E69F00   5. cPurple #CC79A7
%   cGrey #6E6E6E ΔΕΝ είναι σειρά δεδομένων — μόνο πλέγμα, άξονες, βοηθητικά.


% ---------------------------------------------------------------------
%  Α1 — ΚΑΜΠΥΛΗ ΦΡΑΓΜΑΤΩΝ
%  Πότε: όταν το μέγεθος είναι ΑΡΙΘΜΗΤΙΚΗ ΣΥΝΕΠΕΙΑ ενός φράγματος και
%  όχι μέτρηση. Πέντε σειρές, άρα υπόμνημα· κάθε σειρά με ΔΙΚΟ της
%  σύμβολο, ώστε η ταυτότητα να μη στηρίζεται στο χρώμα.
% ---------------------------------------------------------------------
\begin{figure}[t]
\centering
\begin{tikzpicture}
\begin{axis}[
  width=0.80\linewidth, height=5.6cm,
  xlabel={πλήθος ανεκτών σφαλμάτων $f$},
  ylabel={απαιτούμενοι κόμβοι $n$},
  xmin=0.7, xmax=5.3, ymin=0, ymax=28,
  xtick={1,2,3,4,5}, ytick={0,5,10,15,20,25},
  grid=major, grid style={cGrey!20, line width=0.3pt},
  axis line style={cGrey!70},
  tick label style={font=\scriptsize}, label style={font=\small},
  legend style={at={(0.5,-0.30)},anchor=north,legend columns=2,
                font=\scriptsize,draw=cGrey!40,fill=white,
                /tikz/every even column/.append style={column sep=8pt}},
  every axis plot/.append style={thick,mark size=2pt},
]
\addplot[cCFT,   mark=*]         coordinates {(1,3)(2,5)(3,7)(4,9)(5,11)};
  \addlegendentry{$2f{+}1$ — κατάρρευση}
\addplot[cBFT,   mark=square*]   coordinates {(1,4)(2,7)(3,10)(4,13)(5,16)};
  \addlegendentry{$3f{+}1$ — βυζαντινό}
\addplot[cAux,   mark=triangle*] coordinates {(1,5)(2,9)(3,13)(4,17)(5,21)};
  \addlegendentry{$4f{+}1$ — δίκαιη διάταξη}
\addplot[cSand,  mark=diamond*]  coordinates {(1,6)(2,11)(3,16)(4,21)(5,26)};
  \addlegendentry{$5f{+}1$ — μία φάση λιγότερη}
\addplot[cPurple,dashed,mark=o]  coordinates {(1,4)(2,9)(3,14)(4,19)(5,24)};
  \addlegendentry{$5f{-}1$ — κάτω φράγμα δύο βημάτων}
\end{axis}
\end{tikzpicture}
\caption[Πρότυπο Α1: αριθμητική φραγμάτων]%
{\textbf{Πρότυπο Α1.} Καθαρή αριθμητική, όχι μετρήσεις — και η λεζάντα το
λέει πρώτη, γιατί αλλιώς ο αναγνώστης νομίζει ότι βλέπει δεδομένα. Κάθε
σειρά έχει δικό της σύμβολο, ώστε να ξεχωρίζει και σε ασπρόμαυρη εκτύπωση.
Το $5f{-}1$ είναι διακεκομμένο επειδή είναι \emph{κάτω φράγμα} και όχι
κατασκευή: δεύτερη κωδικοποίηση για διαφορετικό είδος μεγέθους.}
\label{fig:protypoA1}
\end{figure}


% ---------------------------------------------------------------------
%  Α2 — ΣΥΓΚΡΙΤΙΚΕΣ ΡΑΒΔΟΙ
%  Πότε: ΜΙΑ μετρική, πολλά συστήματα. Ποτέ δύο άξονες y.
%  Οι τιμές τυπώνονται πάνω από κάθε ράβδο: το #E69F00 και το #CC79A7
%  έχουν αντίθεση κάτω από 3:1 σε λευκό, άρα η ταυτότητα ΔΕΝ επιτρέπεται
%  να στηρίζεται στο χρώμα.
% ---------------------------------------------------------------------
\begin{figure}[t]
\centering
\begin{tikzpicture}
\begin{axis}[
  width=0.80\linewidth, height=5.2cm,
  ybar, bar width=13pt, ymin=0, ymax=78,
  ylabel={μείωση \en{throughput} (\%)},
  symbolic x coords={PBFT,FaB,SBFT,HotStuff,PoE,FTB,Kauri},
  xtick=data, enlarge x limits=0.10,
  x tick label style={rotate=30,anchor=east,font=\tiny},
  tick label style={font=\scriptsize}, label style={font=\small},
  grid=major, ymajorgrids=true, xmajorgrids=false,
  grid style={cGrey!20, line width=0.3pt},
  axis line style={cGrey!70},
  nodes near coords, nodes near coords style={font=\tiny,text=cGrey},
]
\addplot[fill=cBFT!70,draw=cBFT,line width=0.6pt] coordinates
  {(PBFT,65) (FaB,63) (SBFT,55) (HotStuff,45) (PoE,39) (FTB,32) (Kauri,31)};
\end{axis}
\end{tikzpicture}
\caption[Πρότυπο Α2: μία μετρική, πολλά συστήματα]%
{\textbf{Πρότυπο Α2.} Μία σειρά, άρα \emph{κανένα υπόμνημα} — ο τίτλος του
άξονα λέει τι μετριέται. Οι τιμές τυπώνονται πάνω από κάθε ράβδο, ώστε ο
πίνακας να είναι αναγνώσιμος χωρίς να μετρά κανείς με το μάτι. Οριζόντιο
πλέγμα μόνο: το κατακόρυφο δεν προσθέτει τίποτα σε κατηγορικό άξονα.
Η λεζάντα κάθε πραγματικού σχήματος οφείλει να δηλώνει υλικό, μοντέλο
φόρτου, μέγεθος συστάδας και πλήθος επαναλήψεων.}
\label{fig:protypoA2}
\end{figure}


% ---------------------------------------------------------------------
%  Β2 — ΑΚΟΛΟΥΘΙΑ ΜΗΝΥΜΑΤΩΝ  (η TikZ εκδοχή του Mermaid sequenceDiagram)
%  Πότε: όταν σημασία έχει η ΣΕΙΡΑ και το ΠΛΗΘΟΣ των διαδρομών.
%  Κανόνας ευθυγράμμισης: οι κατακόρυφοι άξονες σε ΣΤΑΘΕΡΟ βήμα, τα
%  μηνύματα σε ΣΤΑΘΕΡΟ βήμα ύψους. Καμία διαγώνιος δεν διασταυρώνεται.
% ---------------------------------------------------------------------
\begin{figure}[t]
\centering
\begin{tikzpicture}[
    font=\footnotesize,
    x=3.05cm, y=-0.78cm,                 % y αρνητικό: ο χρόνος πάει κάτω
    msg/.style={->,semithick,cGrey,>={Stealth[length=4pt]}},
    ret/.style={->,semithick,cGrey,dashed,>={Stealth[length=4pt]}},
    role/.style={draw=cCFT,fill=cCFT!8,rounded corners=2pt,
                 minimum width=2.3cm,minimum height=6mm,inner sep=2pt},
    lbl/.style={fill=white,inner sep=1.5pt,font=\scriptsize},
  ]
  % --- ρόλοι, σε σταθερό βήμα x ---
  \foreach \i/\name in {0/πελάτης, 1/ηγέτης, 2/{ακόλουθος 1}, 3/{ακόλουθος 2}}{
    \node[role] (r\i) at (\i,0) {\name};
    \draw[cGrey!55,line width=0.5pt] (\i,0.45) -- (\i,7.6);
  }
  % --- μηνύματα, σε σταθερό βήμα y ---
  \draw[msg] (0,1.4) -- node[lbl,above]{αίτημα}            (1,1.4);
  \draw[msg] (1,2.4) -- node[lbl,above]{πρόταση}           (2,2.4);
  \draw[msg] (1,3.4) -- node[lbl,above]{πρόταση}           (3,3.4);
  \draw[ret] (2,4.4) -- node[lbl,above]{ψήφος}             (1,4.4);
  \draw[ret] (3,5.4) -- node[lbl,above]{ψήφος}             (1,5.4);
  \node[draw=cSand,fill=cSand!12,rounded corners=2pt,inner sep=3pt,
        font=\scriptsize,anchor=west] at (0.62,6.4)
       {πλειοψηφία: η εντολή δεσμεύεται};
  \draw[msg] (1,7.3) -- node[lbl,above]{απόκριση}          (0,7.3);
\end{tikzpicture}
\caption[Πρότυπο Β2: ακολουθία μηνυμάτων]%
{\textbf{Πρότυπο Β2.} Ο χρόνος τρέχει προς τα κάτω και κάθε μήνυμα παίρνει
δική του γραμμή: έτσι το πλήθος των διαδρομών \emph{μετριέται} αντί να
εκτιμάται. Συνεχής γραμμή = πρόταση του ηγέτη, διακεκομμένη = απόκριση
ακολούθου — η κατεύθυνση φαίνεται και χωρίς χρώμα. Το πλαίσιο σημειώνει τη
στιγμή που η εντολή γίνεται οριστική, που είναι το μόνο σημείο του
σχήματος με σημασία για την ορθότητα.}
\label{fig:protypoB2}
\end{figure}
